% Converted from README.md
\subsection{Java OOP primer - Game of Life}\label{java-oop-primer---game-of-life}

Program predstavlja implementaciju poznate ćelijske automatizacije pod imenom \href{https://en.wikipedia.org/wiki/Conway\%27s_Game_of_Life}{Game of Life}. Početna konfiguracija igre (početno stanje table kao i veličina mreže) su konfigurabilni. Stanje igre se ispisuje na standardni izlaz nakon svake iteracije.

Da bismo pisali testove jedinica koda, moramo razumeti odnose klasa u okviru aplikacije:

\begin{itemize}
\tightlist
\item
  paket \texttt{app} sadrži implementaciju \texttt{GameOfLife} igre
\item
  paket \texttt{model} sadrži implementacije klasa \texttt{Cell} i \texttt{Grid}, koje se koriste za reprezentaciju stanja igre

  \begin{itemize}
  \tightlist
  \item
    paket \texttt{model.conf} sadrži implementacije početnih stanja igre (nasumičnu i jednu unapred kreiranu sa specifičnim osobinama, pod nazivom \texttt{Glider})
  \end{itemize}
\end{itemize}

Testovi koje bismo voleli da napišemo bi testirali:

\begin{itemize}
\tightlist
\item
  da li se mreža ispravno kreira na osnovu konfiguracije
\item
  da li se pravila igre ispravno primenjuju iz generacije u generaciju
\item
  da li se stanje igre ispravno ispisuje na standardni izlaz
\end{itemize}

Testove da li se mreža ispravno kreira na osnovu konfiguracije možemo lako napraviti. \texttt{Grid} klasa već implementira sve neophodno. Jedini problem predstavlja to što nemamo način da ručno postavimo konfiguraciju koja će se koristiti kao početna. Možemo testirati nasumičnu konfiguraciju međutim radije bismo da imamo deterministične testove. Takođe, naši testovi bi idealno testirali konfiguraciju u specijalnim slučajevima, npr. ivice mreže, što ne možemo kontrolisati nasumičnom konfiguracijom. Možemo ručno promeniti stanje klase \texttt{Grid} nakon kreiranja ali testovi \textbf{nikada} ne treba da zalaze u detalje implementacije klasa, tj. njihova stanja. Kreiranje ručnih determinističkih konfiguracija nam omogućava i testiranje pravila igre - možemo kreirati \texttt{Grid} objekat sa specifičnom konfiguracijom namenjenom da testira određeno pravilo igre ili kombinaciju više pravila, pozivajući metod \texttt{Grid.advance()} i testirati da li su se ćelije promenile onako kako bi trebalo. Zatim možemo testirati i specijalne slučajeve kao što su ivice mreže kao i situacije u kojima se na jednu ćeliju primenjuje više pravila.

Da bismo implementirali ručne konfiguracije, primetimo da je \texttt{model.conf.GliderConfiguration} jedna implementacija ručne konfiguracije. Možemo da implementiramo naše specifične konfiguracije na sličan način, međutim imali bismo ponavljanje koda pošto bi se jedino menjao konstruktor klase. Apstrahujmo \texttt{GliderConfiguration} implementaciju - dodajemo klasu \texttt{ManualConfiguration} i menjamo klasu \texttt{GliderConfiguration} da nasleđuje klasu \texttt{ManualConfiguration}. Sada naše test konfiguracije mogu da instanciraju \texttt{ManualConfiguration} sa odgovarajućom konfiguracijom mreže, a takođe aplikacija može da se proširi dodavanjem novih predefinisanih konfiguracija.

Slično kao u \texttt{02\_refactoring/01\_cpp\_calculator} primeru, treba apstrahovati rad sa standardnim izlazom kako bismo mogli da testiramo prikaz stanja igre. Prikaz stanja se trenutno vrši u glavnoj klasi aplikacije, što takođe nije optimalno. Dodajmo \texttt{views} paket sa implementacijom \texttt{View} interfejsa koji predstavlja apstraktnu implementaciju prikaza aplikacije. Sada možemo kreirati implementaciju koja ispisuje stanje igre na proizvoljni izlazni tok (\texttt{PrintStreamView}) odnosno \texttt{System.out} ukoliko izlazni tok nije naveden. Logiku ispisa stanja igre pomeramo iz \texttt{app.GameOfLife} u \texttt{view.ConsoleView}.
